\chapter{Maxima, Minima, and the Shape of a Graph}

Here is a question so ordinary you may never have thought to ask it carefully:
how do you \emph{know} what a graph looks like? You have surely sketched curves
before --- plotted a few points, connected the dots, trusted that the picture
between the dots behaved itself. But that is guessing. A function can do
surprising things between two sample points. It can dive down and back up. It can
have a sharp bend you would never see at the resolution of your dots. To draw a
graph \emph{correctly} --- to know where it rises, where it falls, where it turns
around, where it bulges up and where it sags down --- you need more than a handful
of values. You need the derivative, and you need the derivative of the derivative.

Two derivatives are enough to pin the picture down. The first
derivative $f'$ tells you the \emph{direction} the curve is heading --- uphill or
downhill. The second derivative $f''$ tells you how that direction is
\emph{changing} --- whether the curve bends upward or downward. Put
the two together and you can determine the overall shape of a graph,
and find its highest and lowest points.
Those highest and lowest points are the payoff. They are what optimization is made
of: the cheapest design, the strongest beam, the fastest route, the maximum
profit, the least error. All of it begins here, with the simple observation that
\emph{at the top of a hill, the ground is momentarily level}.

Everything in this chapter rests on the Mean Value Theorem from the previous
chapter. The MVT was the bridge from $f'$ back to $f$. Now we cross that bridge
over and over. When we claim ``$f'>0$ means $f$ is increasing,'' that is the MVT
talking. When we claim a sign change in $f'$ means a genuine peak, that is the MVT
again. We will keep pointing back to it, because the real reason these tests
\emph{work} --- and are not just plausible folklore --- is that one theorem.

\section{The top of the hill is level}

Stand at the very top of a smooth hill. Look at the ground right at your feet.
It is flat. Not flat far away --- the hill falls off on both sides --- but flat
\emph{right there}, at the summit. Geometry forces this. If the ground at the summit were tilted even slightly uphill in some
direction, then you could walk that way and go higher --- so you were not at the
summit after all. The only way to be at the very top is for every direction to be
level or downhill, and for a smooth curve in one variable that means the tangent
line is horizontal.

\begin{intuitionbox}
At a peak or a valley of a smooth graph, the tangent line is horizontal: the
slope is zero. The reason is almost a tautology once you see it. To the left of a
peak the function is climbing ($f'>0$); to the right it is falling ($f'<0$). A
continuous slope cannot jump from positive to negative without passing through
zero. So right at the top, $f'=0$. The summit is where ``going up'' turns into
``going down,'' and the turning point is level.
\end{intuitionbox}

Let us make the words precise, because ``highest point'' hides a real subtlety:
highest compared to \emph{what}? Compared to the entire domain, or just compared
to nearby points?

\begin{definition}[Local and global extrema]\label{ch16:def-extrema}
Let $f$ be defined on a set $D$ and let $c \in D$.
\begin{itemize}
\item $f$ has a \emph{local maximum} at $c$ if $f(c) \ge f(x)$ for all $x$ in some
open interval around $c$ (a local minimum if $f(c)\le f(x)$ nearby).
\item $f$ has a \emph{global} (or \emph{absolute}) \emph{maximum} at $c$ if
$f(c)\ge f(x)$ for \emph{all} $x \in D$ (global minimum if $f(c)\le f(x)$ for all
$x\in D$).
\end{itemize}
A \emph{local extremum} (plural: \emph{extrema}, from the Latin for outermost) is either a local max or a local min.
\end{definition}

A local maximum is a peak relative to its immediate surroundings --- the top of
\emph{a} hill, even if a taller mountain stands a mile away. A global maximum is
the highest point across the entire domain of the function. Every interior global extremum is automatically a
local one (being highest overall makes you highest locally), but not the reverse.

\begin{figure}[H]
\centering
\begin{tikzpicture}
\begin{axis}[
  width=12cm, height=6.2cm,
  axis lines=middle,
  xlabel={$x$}, ylabel={$y$},
  xmin=-0.4, xmax=4.6, ymin=-1.4, ymax=6.2,
  xtick=\empty, ytick=\empty,
  domain=0:4.2, samples=120,
  clip=false,
]
\addplot[curveblue, very thick] {x^3 - 6*x^2 + 9*x} ;
\addplot[only marks, mark=*, mark size=1.6pt, tangentred]
  coordinates {(1,4)};
\addplot[only marks, mark=*, mark size=1.6pt, tangentred]
  coordinates {(3,0)};
\node[anchor=south, tangentred] at (axis cs:1,4.15) {local max};
\node[anchor=north, tangentred] at (axis cs:3,-0.25) {local min};
\addplot[tangentred!80!black, thick, dashed, domain=0.35:1.65] {4};
\addplot[tangentred!80!black, thick, dashed, domain=2.35:3.65] {0};
\node[anchor=west, curveblue] at (axis cs:3.9,5.4) {$f$};
\end{axis}
\end{tikzpicture}
\caption{At a local maximum and a local minimum of a smooth curve, the tangent
line (dashed) is horizontal. The peak is where climbing turns into falling.}
\label{ch16:fig-extrema}
\end{figure}

\subsection{Fermat's theorem: where to look}

The flatness-at-the-summit idea has a name and a proof. It is one of the oldest
results in the subject, found by Pierre de Fermat decades before Newton and
Leibniz had a systematic calculus.

\begin{theorem}[Fermat's theorem on extrema]\label{ch16:thm-fermat}
If $f$ has a local extremum at $c$, and $f$ is differentiable at $c$, then
$f'(c)=0$.
\end{theorem}

\begin{intuitionbox}
The summit argument translates directly into algebra, exposing why
differentiability is the essential hypothesis. The derivative is a single number,
the limit of the difference quotients $\dfrac{f(c+h)-f(c)}{h}$ as $h\to 0$.
Approaching from the right and from the left, both one-sided limits must converge
to that same value. At a local maximum, the numerator $f(c+h)-f(c)$ is $\le 0$
for all sufficiently small $h$ (no nearby point is higher).
When $h>0$, the quotient is nonpositive divided by positive, so $\le 0$.
When $h<0$, the quotient is nonpositive divided by negative, so $\ge 0$.
The right-hand limit forces $f'(c) \le 0$, while the left-hand limit forces
$f'(c) \ge 0$. The only real number that is simultaneously $\le 0$ and $\ge 0$ is
$0$.
\end{intuitionbox}

\begin{proof}
Suppose $f$ has a local maximum at $c$ (the minimum case is identical with
inequalities flipped). Then there is an interval around $c$ on which
$f(x)\le f(c)$, i.e.\ $f(c+h)-f(c)\le 0$ for all small $h$.

For $h>0$,
\[
\frac{f(c+h)-f(c)}{h} \le 0
\qquad\text{(nonpositive over positive)},
\]
so the right-hand limit satisfies $f'(c)=\lim_{h\to 0^+}\dfrac{f(c+h)-f(c)}{h}\le
0$.

For $h<0$,
\[
\frac{f(c+h)-f(c)}{h} \ge 0
\qquad\text{(nonpositive over negative)},
\]
so the left-hand limit satisfies $f'(c)=\lim_{h\to 0^-}\dfrac{f(c+h)-f(c)}{h}\ge
0$.

Because $f$ is differentiable at $c$, these two one-sided limits are equal and both
equal $f'(c)$. The only way to have $f'(c)\le 0$ and $f'(c)\ge 0$ at once is
$f'(c)=0$.
\end{proof}

\begin{historybox}
The discovery that extrema occur where rates of change vanish predates both Newton and Leibniz. In 1636, Pierre de Fermat circulated a short manuscript, \emph{Methodus ad disquirendam maximam et minimam}, presenting his method of \emph{adequality} (from the Latin for ``approximate equality''). Fermat was inspired by an observation of the ancient Greek geometer Pappus of Alexandria on dividing a line segment: at a maximum, a slight nudge in the input produces virtually no change in the output. Fermat replaced $x$ with $x + e$, set $f(x+e)$ ``adequal'' to $f(x)$, divided through by the perturbation $e$, and then suppressed all terms still containing $e$. In modern language, Fermat was evaluating $\lim_{e\to 0}\frac{f(x+e)-f(x)}{e} = 0$, identifying peaks and troughs by finding stationary values forty years before the invention of calculus.

This procedure immediately triggered a fierce controversy. In 1638, Ren\'e Descartes sharply criticized Fermat's method, demanding to know how Fermat could legitimately treat $e$ as a nonzero quantity to justify dividing by it, and then instantly discard it as zero without creating a logical contradiction. Without a rigorous definition of a limit, Fermat could only plead that the method consistently produced the correct geometric answers.

The second half of the story---the geometry of bending and concavity---belongs to Gottfried Wilhelm Leibniz. In his 1684 paper and subsequent memoirs in \emph{Acta Eruditorum}, Leibniz analyzed what he called \emph{puncta flexus contrarii}: inflection points where a curve transitions between bending upward and bending downward. Leibniz realized that just as the first differential $\dd y$ determines the slope, the second differential $\dd^2 y$ determines the bending or curvature. At an inflection point, $\dd^2 y = 0$, and the tangent line actually slices through the curve rather than merely grazing one side.

Fermat and Leibniz understood the geometric intuition perfectly, but lacked the limit concept to justify why vanishing terms could be cancelled. In the nineteenth century, Cauchy and Weierstrass provided the rigorous basis with the $\eps$-$\delta$ limit and the Mean Value Theorem, proving that Fermat's stationary points and Leibniz's curvature transitions are exact consequences of real analysis. This foundation supported the calculus of variations, Fermat's principle of least time in optics, and the principle of least action in classical physics. Yet the story is far from closed: while single-variable optimization is completely solved, finding global extrema for non-convex functions of millions of variables---such as the objective surfaces traversed when training deep neural networks---is computationally intractable (NP-hard) in general. Understanding high-dimensional geometry, saddle points, and optimization dynamics remains one of the most vibrant open frontiers of modern mathematics.
\end{historybox}

\begin{trapbox}
Fermat's theorem is a one-way street, and reading it backwards is the single most
common mistake in this whole subject. $f'(c)=0$ does \emph{not} guarantee an
extremum. The graph of $f(x)=x^3$ has $f'(0)=0$, yet $0$ is neither a max nor a
min --- the curve flattens for an instant and keeps right on rising. A zero
derivative is a \emph{candidate}, a place worth checking, nothing more. The
theorem tells you where to look, not what you will find.
\end{trapbox}

\subsection{Critical points: candidate extrema}

Fermat's theorem says smooth extrema hide where $f'=0$. But there is a second way
a function can turn around: at a \emph{corner}, where the derivative does not
exist at all. The function $f(x)=\abs{x}$ has its minimum at $0$, and $f'(0)$
simply does not exist --- the slope is $-1$ on the left and $+1$ on the right, with
no single tangent. Fermat's theorem said nothing here, because its hypothesis
(``differentiable at $c$'') failed. So our search must include both kinds of points.

\begin{definition}[Critical point]\label{ch16:def-critical}
A \emph{critical point} of $f$ is a number $c$ in the domain of $f$ at which either
$f'(c)=0$ or $f'(c)$ does not exist. The value $f(c)$ is a \emph{critical value}.
\end{definition}

\begin{focusbox}
Combining Fermat's theorem with the corner case gives the rule that organizes the
entire hunt for extrema:
\[
\boxed{\text{Every local extremum of } f \text{ occurs at a critical point.}}
\]
This does not say critical points are extrema --- most are not. It says you never
have to search the whole real line. The peaks and valleys can \emph{only} sit at
the finitely many places where $f'$ is zero or undefined. Find those, and you have
narrowed the candidate extrema to a finite set. That is the entire strategy.
\end{focusbox}

\begin{example}[Listing critical points]\label{ch16:ex-list-crit}
Find the critical points of $f(x)=x^3-3x^2-9x+5$.

\smallskip
Differentiate: $f'(x)=3x^2-6x-9=3(x^2-2x-3)=3(x-3)(x+1)$. This is a polynomial,
defined everywhere, so the only critical points come from $f'=0$:
\[
3(x-3)(x+1)=0 \;\Longrightarrow\; x=3 \ \text{ or } \ x=-1.
\]
Two suspects: $x=-1$ and $x=3$. We do not yet know which (if either) is a max or a
min --- that is the next section's job.
\end{example}

\begin{example}[A critical point from a missing derivative]\label{ch16:ex-cusp}
Find the critical points of $g(x)=x^{2/3}$.

\smallskip
Here $g'(x)=\dfrac{2}{3}x^{-1/3}=\dfrac{2}{3\,x^{1/3}}$. There is no $x$ making
this zero (the numerator is a constant $2/3$). But at $x=0$ the formula blows up:
$g'(0)$ does not exist. Yet $0$ is in the domain of $g$, and $g(0)=0$ while
$g(x)>0$ for every other $x$. So $x=0$ is a critical point --- in fact a global
minimum --- found \emph{only} because we remembered to check where the derivative
fails to exist. The graph has a sharp \emph{cusp} (a sharp point where two branches meet with a vertical tangent line) pointing down at the origin.
\end{example}

\begin{figure}[H]
\centering
\begin{tikzpicture}
\begin{axis}[
  name=axCorner, title style={font=\small, yshift=4pt},
  width=6.3cm, height=5.2cm,
  axis lines=middle, title={corner: $f(x)=|x|$},
  xlabel={$x$}, ylabel={$y$},
  xmin=-2.4, xmax=2.4, ymin=-0.4, ymax=2.5,
  xtick={-2,2}, ytick={1},
  domain=-2.1:2.1, samples=60, clip=false,
]
\addplot[curveblue, very thick] {abs(x)};
\addplot[only marks, mark=*, mark size=2pt, tangentred] coordinates {(0,0)};
\node[anchor=north east, tangentred] at (axis cs:-0.08,-0.05) {\footnotesize min $(0,0)$};
\node[anchor=north east, gray!60!black] at (axis cs:-0.7,0.7) {\footnotesize slope $-1$};
\node[anchor=north west, gray!60!black] at (axis cs:0.7,0.7) {\footnotesize slope $+1$};
\node[anchor=south east, trapred] at (axis cs:-0.15,1.7) {\footnotesize $f'(0)$ undefined};
\end{axis}
\begin{axis}[
  at={(axCorner.east)}, anchor=west, xshift=1.2cm,
  title style={font=\small, yshift=4pt},
  width=6.3cm, height=5.2cm,
  axis lines=middle, title={cusp: $g(x)=x^{2/3}$},
  xlabel={$x$}, ylabel={$y$},
  xmin=-2.4, xmax=2.4, ymin=-0.4, ymax=2.5,
  xtick={-2,2}, ytick={1},
  clip=false,
]
\addplot[curveblue, very thick, domain=-1.25:1.25, samples=120] ({x^3}, {x^2});
\addplot[only marks, mark=*, mark size=2pt, tangentred] coordinates {(0,0)};
\node[anchor=north east, tangentred] at (axis cs:-0.08,-0.05) {\footnotesize min $(0,0)$};
\node[anchor=south west, trapred] at (axis cs:0.15,1.7) {\footnotesize vertical tangent};
\end{axis}
\end{tikzpicture}
\caption{Two critical points where the derivative does not exist. Left: a corner,
where left- and right-hand slopes disagree ($-1$ versus $+1$). Right: a cusp,
where the tangent line turns vertical and the slope blows up. In both cases,
$x=0$ is a genuine minimum that Fermat's equation $f'=0$ would overlook.}
\label{ch16:fig-corner-cusp}
\end{figure}

\begin{recoverybox}
\textbf{Pause and think}: Suppose a function $f$ is continuous everywhere, climbs steadily on $(-\infty, 2)$, and falls steadily on $(2, \infty)$. Does $x=2$ have to be a local maximum? Yes: by definition, $f(2) \ge f(x)$ for all nearby $x$. But does $f'(2)$ have to equal $0$? Not at all: the crest could be a sharp corner like $f(x) = 3 - |x-2|$, where the two one-sided slopes disagree and the derivative does not exist. That is why Fermat's equation $f'=0$ is only half the net: critical points must include every point where $f'$ fails to exist.
\end{recoverybox}

\section{The first-derivative test: reading the sign of the slope}

We have a list of candidate points. Now we classify them. The tool is a direct
geometric principle: \emph{the sign of $f'$ tells you which way the
function is going.} Positive slope, going up. Negative slope, going down. So to
decide whether a critical point is a peak, a valley, or neither, we just watch what
the slope does as we walk through it.

\begin{whybox}
Why does the sign of $f'$ control whether $f$ rises or falls? This is the
Monotonicity Theorem (Theorem~\ref{ch15:thm-mono} in the preceding chapter),
resting squarely on the Mean Value Theorem. For any two points $a<b$ in an interval, the MVT guarantees
a point $\xi \in (a,b)$ satisfying
\[
f(b)-f(a)=f'(\xi)\,(b-a).
\]
Because $b>a$, the factor $(b-a)$ is strictly positive. Consequently, the sign of the
vertical increment $f(b)-f(a)$ matches the sign of the instantaneous slope $f'(\xi)$.
If $f'>0$ across the entire interval, then $f'(\xi)>0$ for every pair of sample points,
forcing $f(b)>f(a)$ throughout: $f$ is strictly increasing. If $f'<0$ throughout,
$f(b)<f(a)$, and $f$ is strictly decreasing. The first-derivative test is the direct
application of this relationship around a critical point.
\end{whybox}

\begin{theorem}[First-derivative test]\label{ch16:thm-fdt}
Suppose $f$ is continuous at a critical point $c$ and differentiable on an open
interval around $c$ (except possibly at $c$ itself). Consider the sign of $f'$ just
to the left and just to the right of $c$:
\begin{itemize}
\item if $f'$ changes from $+$ to $-$ at $c$, then $f$ has a \emph{local maximum}
at $c$;
\item if $f'$ changes from $-$ to $+$ at $c$, then $f$ has a \emph{local minimum}
at $c$;
\item if $f'$ does \emph{not} change sign at $c$ (stays $+$ or stays $-$), then $c$
is \emph{not} a local extremum.
\end{itemize}
\end{theorem}

\begin{proof}
Suppose $f'>0$ on an interval $(a,c)$ to the left of $c$ and $f'<0$ on an interval
$(c,b)$ to the right. By the Monotonicity Theorem, $f$ is increasing on $(a,c)$ and
decreasing on $(c,b)$; by continuity at $c$ these extend to $f$ increasing on
$(a,c]$ and decreasing on $[c,b)$. So for any $x$ in $(a,c)$ we have $f(x)<f(c)$,
and for any $x$ in $(c,b)$ we have $f(x)<f(c)$. Thus $f(c)\ge f(x)$ for all $x$ in
$(a,b)$: a local maximum. The minimum case reverses the inequalities. If $f'$ keeps
the same sign on both sides, $f$ is monotone straight through $c$ --- still rising
(or still falling) --- so $c$ cannot be a turning point.
\end{proof}

\begin{intuitionbox}
Picture a ball rolling along the graph. Where $f'>0$ the ground slopes up to the
right, so a ball would roll \emph{left}, downhill. Where $f'<0$ the ground slopes
down to the right, so the ball rolls \emph{right}. A local \emph{minimum} is a spot
both neighbors roll \emph{toward} (down from the left, down from the right) --- a
basin, $-$ then $+$. A local \emph{maximum} is a spot both neighbors roll
\emph{away} from --- a ridge, $+$ then $-$. And if the ball rolls the same
direction on both sides, it just passes through; no resting place.
\end{intuitionbox}

The practical method is a \emph{sign chart}. Mark the critical points on a number
line; they cut it into intervals. On each interval, $f'$ has a constant sign
(it can only change sign by passing through a zero or a point of nondefinition ---
a critical point --- so between consecutive critical points the sign is locked).
Test one convenient point per interval and record $+$ or $-$.

\begin{example}[First-derivative test in full]\label{ch16:ex-fdt}
Classify the critical points of $f(x)=x^3-3x^2-9x+5$ from
Example~\ref{ch16:ex-list-crit}.

\smallskip
We had $f'(x)=3(x-3)(x+1)$, with critical points $x=-1$ and $x=3$. These split the
line into three intervals. Pick a test point in each and read the sign of $f'$ from
the factored form:
\[
\begin{array}{c|ccc}
\text{interval} & (-\infty,-1) & (-1,3) & (3,\infty)\\\hline
\text{test } x & -2 & 0 & 4\\
(x-3) & - & - & +\\
(x+1) & - & + & +\\
f'=3(x-3)(x+1) & + & - & +\\
f & \nearrow & \searrow & \nearrow
\end{array}
\]
At $x=-1$, $f'$ goes $+\to-$: a \textbf{local maximum}, value
$f(-1)=-1-3+9+5=10$. At $x=3$, $f'$ goes $-\to+$: a \textbf{local minimum}, value
$f(3)=27-27-27+5=-22$. The function climbs to a peak at $(-1,10)$, slides down into
a valley at $(3,-22)$, then climbs forever.
\end{example}

\begin{figure}[H]
\centering
\begin{tikzpicture}
\begin{axis}[
  width=12cm, height=6.4cm,
  axis lines=middle,
  xlabel={$x$}, ylabel={$y$},
  xmin=-3.2, xmax=5.4, ymin=-26, ymax=18,
  xtick={-1,3}, ytick=\empty,
  domain=-2.6:5, samples=160,
  clip=false,
]
\addplot[curveblue, very thick] {x^3 - 3*x^2 - 9*x + 5};
\addplot[only marks, mark=*, mark size=1.6pt, tangentred]
  coordinates {(-1,10) (3,-22)};
\node[anchor=south, tangentred] at (axis cs:-1,11) {$(-1,10)$ max};
\node[anchor=north, tangentred] at (axis cs:3,-23.5) {$(3,-22)$ min};
\node[green!50!black] at (axis cs:-2.0,-14) {$f'>0$};
\node[trapred] at (axis cs:1.0,-14) {$f'<0$};
\node[green!50!black] at (axis cs:4.3,-14) {$f'>0$};
\end{axis}
\end{tikzpicture}
\caption{$f(x)=x^3-3x^2-9x+5$. The slope is positive, then negative, then positive
--- so the curve rises, falls, rises, with a peak at $x=-1$ and a valley at $x=3$.}
\label{ch16:fig-cubic}
\end{figure}

\begin{example}[A critical point that is not an extremum]\label{ch16:ex-no-extremum}
Classify the critical point of $f(x)=x^3$.

\smallskip
$f'(x)=3x^2$, zero only at $x=0$. But $3x^2\ge 0$ everywhere: $f'>0$ for $x<0$
\emph{and} $f'>0$ for $x>0$. The slope does not change sign --- it dips to zero and
bounces back up. By the first-derivative test, $x=0$ is \emph{not} an extremum. The
curve flattens for one instant and keeps climbing. This is the trap from earlier
made concrete: $f'(0)=0$ did not deliver an extremum.
\end{example}

\section{Concavity: which way the curve bends}

Knowing where a graph rises and falls still leaves its \emph{shape} ambiguous. A
curve can climb from $A$ to $B$ in two very different ways: bulging up like the top
of a dome, or sagging down like the inside of a bowl. Both go up; they have
opposite \emph{curvature}. To capture that, we ask not how $f$ is changing, but how
$f'$ is changing --- and the rate of change of $f'$ is $f''$.

\begin{definition}[Concavity]\label{ch16:def-concave}
On an interval where $f$ is differentiable:
\begin{itemize}
\item $f$ is \emph{concave up} if $f'$ is increasing on the interval (the slope
gets ever larger --- more positive or less negative as you move right);
\item $f$ is \emph{concave down} if $f'$ is decreasing on the interval.
\end{itemize}
\end{definition}

Apply the Monotonicity Theorem to the function $f'$ itself: $f'$ is increasing
exactly where \emph{its} derivative $f''$ is positive. That gives the working test.

\begin{focusbox}
\[
f''>0 \;\Longrightarrow\; \text{concave up (cup-shaped, }\smile\text{)},
\qquad
f''<0 \;\Longrightarrow\; \text{concave down (cap-shaped, }\frown\text{)}.
\]
A memory hook: $f''>0$ holds \emph{water}, like a cup. $f''<0$ spills it, like a
frown.
\end{focusbox}

\begin{intuitionbox}
There are three ways to picture concavity, and you should hold all three at
once.
\begin{enumerate}
\item \textbf{The slope is turning.} Concave up means the tangent line is rotating
counterclockwise as you move right --- the slope keeps increasing. Concave down
means it rotates clockwise.
\item \textbf{The curve vs.\ its tangent.} A concave-up curve lies \emph{above}
each of its tangent lines (the tangent underestimates it). A concave-down curve
lies \emph{below} its tangents. This is why the linear approximation from the
tangent-line chapter consistently undershoots a convex function and overshoots a
concave one.
\item \textbf{The chord.} A concave-up curve lies \emph{below} the straight chord
joining two of its points (the straight line segment connecting two points on the
curve; think of the bottom of a valley); concave down lies above its chords. This is
the geometric heart of convexity, and it is what makes inequalities like the AM--GM
inequality fall out of a single picture.
\end{enumerate}
\end{intuitionbox}

\begin{figure}[H]
\centering
\begin{tikzpicture}
\begin{axis}[
  name=ax1, title style={font=\small, yshift=4pt},
  width=6.4cm, height=5cm,
  axis lines=middle, title={concave up: $f''>0$},
  xlabel={$x$}, xmin=-2.3, xmax=2.3, ymin=-0.4, ymax=4.4,
  xtick=\empty, ytick=\empty, domain=-2:2, samples=80, clip=false,
]
\addplot[curveblue, very thick] {x^2};
\addplot[tangentred, thick, domain=-2:0.3] {-2*x-1};
\node[curveblue, anchor=west] at (axis cs:1.1,3.6) {$f$};
\node[tangentred, font=\footnotesize, anchor=north] at (axis cs:-1.4,1.2) {tangent below};
\end{axis}
\begin{axis}[
  at={(ax1.east)}, anchor=west, xshift=1.2cm, title style={font=\small, yshift=4pt},
  width=6.4cm, height=5cm,
  axis lines=middle, title={concave down: $f''<0$},
  xlabel={$x$}, xmin=-2.3, xmax=2.3, ymin=-4.4, ymax=0.6,
  xtick=\empty, ytick=\empty, domain=-2:2, samples=80, clip=false,
]
\addplot[curveblue, very thick] {-x^2};
\addplot[tangentred, thick, domain=-0.3:2] {1-2*x};
\node[curveblue, anchor=east] at (axis cs:-1.1,-2.0) {$f$};
\node[tangentred, font=\footnotesize, anchor=west] at (axis cs:1.1,-0.8) {tangent above};
\end{axis}
\end{tikzpicture}
\caption{Concave up ($\smile$): the curve sits above its tangent lines. Concave
down ($\frown$): the curve sits below its tangent lines.}
\label{ch16:fig-concavity}
\end{figure}

\begin{whybox}
Why does a concave-up curve ($f''>0$) lie above its tangent lines?
Fix a point $a$, and let
\[
T(x)=f(a)+f'(a)(x-a)
\]
be the tangent line at $a$. Define the gap $g(x)=f(x)-T(x)$, measuring how far the
curve sits above the tangent line. Note that $g(a)=0$ and $g'(x)=f'(x)-f'(a)$, so
$g'(a)=0$.

When $f''>0$ on the interval, $f'$ is strictly increasing:
\begin{itemize}
\item For $x>a$: $f'(x)>f'(a)$, so $g'(x)>0$. Since $g(a)=0$ and $g$ is increasing
to the right of $a$, we have $g(x)>0$.
\item For $x<a$: $f'(x)<f'(a)$, so $g'(x)<0$. Since $g$ is decreasing toward $a$
where $g(a)=0$, moving leftward from $a$ means $g(x)>0$.
\end{itemize}
On either side of $a$, $g(x)\ge 0$: the curve sits strictly above its tangent line
except at the point of tangency itself. The entire deduction is two successive
applications of the Mean Value Theorem: first showing $f''>0$ makes $f'$ increase,
then showing $g'>0$ makes $g$ increase.
\end{whybox}

\begin{recoverybox}
\textbf{Pause and think}: Could a curve be concave up everywhere ($f''>0$ on all of $\R$) and still possess a local maximum? Picture what that would demand: at a local maximum the tangent line is horizontal, and concave up forces the curve to lie \emph{above} that tangent line on both sides. But lying above a horizontal tangent means climbing even higher, which destroys the peak! Therefore, a function that is strictly concave up everywhere can never have a local maximum.
\end{recoverybox}

\subsection{Inflection points: where the bend reverses}

A curve does not have to keep the same concavity forever. Many of the most
important shapes in science change their bending partway through. An inflection
point is where that switch happens.

\begin{definition}[Inflection point]\label{ch16:def-inflection}
A point $\big(c,f(c)\big)$ on the graph of a continuous function is an
\emph{inflection point} if the concavity of $f$ \emph{changes} there --- from up to
down, or down to up.
\end{definition}

Since concavity is governed by the sign of $f''$, an inflection point is a place
where $f''$ changes sign. And just as $f'$ can only change sign at a zero or a gap,
$f''$ can only change sign where $f''=0$ or $f''$ fails to exist. Those are the
\emph{candidate} inflection points --- and, exactly as with Fermat's theorem,
candidacy is not guilt.

\begin{trapbox}
$f''(c)=0$ does \emph{not} by itself make $\big(c,f(c)\big)$ an inflection point.
You must check that $f''$ actually \emph{changes sign}. Classic counterexample:
$f(x)=x^4$ has $f''(x)=12x^2$, which is zero at $x=0$ but $\ge 0$ on both sides ---
no sign change, no inflection. The curve is concave up throughout; it merely
flattens momentarily at the bottom. Always confirm the sign flip, never read it off
from $f''=0$ alone.
\end{trapbox}

\begin{example}[The S-curve and its inflection]\label{ch16:ex-logistic}
Consider $f(x)=\dfrac{1}{1+e^{-x}}$, the logistic curve, widely used in
population dynamics, neural networks, and models of saturating growth. Find
where it is concave up, concave down, and its inflection point.

\smallskip
Write $s=f(x)$. Differentiating with the quotient or chain rule yields
$f'(x)=s(1-s)$, and a second derivative gives
\[
f''(x)=f'(x)(1-2s)=s(1-s)(1-2s).
\]
Now $s$ ranges strictly between $0$ and $1$, so $s>0$ and $(1-s)>0$ always; the
sign of $f''$ is determined purely by the factor $(1-2s)$. That is positive when
$s<\tfrac12$ and negative when $s>\tfrac12$. And $s=\tfrac12$ happens exactly at
$x=0$. So:
\[
f''>0 \text{ for } x<0 \ (\text{concave up}),\qquad
f''<0 \text{ for } x>0 \ (\text{concave down}).
\]
The concavity flips at $x=0$: an inflection point at $\big(0,\tfrac12\big)$. This
is the famous ``tipping point'' of the S-curve --- the moment of fastest growth,
where the curve stops accelerating and begins to level off toward saturation.
\end{example}

\begin{applicationbox}
In applications, the inflection point is often \emph{the} quantity of
interest. In an epidemic modeled by a logistic curve, the inflection point is the
day new cases \emph{peak} --- after it, growth is still happening but slowing, and
the worst is mathematically behind you even though the total is still rising. In a
product-adoption curve it is the moment of maximum sales velocity. In both cases
$f$ (the cumulative total) is still increasing the whole time; it is $f''$ changing
sign that signals the turn. The first derivative tells you the situation is getting
worse; the second derivative tells you how fast it is getting worse, and its sign
change is the first good news.
\end{applicationbox}

\begin{figure}[H]
\centering
\begin{tikzpicture}
\begin{axis}[
  width=12cm, height=6.2cm,
  axis lines=middle,
  xlabel={$x$}, ylabel={$y$},
  xmin=-6.2, xmax=6.2, ymin=-0.15, ymax=1.25,
  xtick={0}, ytick=\empty,
  domain=-6:6, samples=160, clip=false,
  title style={font=\small},
]
\addplot[gray, dashed] coordinates {(-6,1) (6,1)};
\node[anchor=south east, gray!70!black] at (axis cs:6,1.02) {\footnotesize $y=1$};
\addplot[curveblue, very thick] {1/(1+exp(-x))};
\addplot[tangentred, thick, dashed, domain=-3.5:3.5] {0.5 + 0.25*x};
\node[tangentred, font=\footnotesize, anchor=west] at (axis cs:2.85, 1.15) {tangent line};
\addplot[only marks, mark=*, mark size=2pt, tangentred] coordinates {(0,0.5)};
\draw[<-, >=stealth, tangentred, shorten <=2pt] (axis cs:0,0.5) -- (axis cs:-0.7,0.65)
  node[anchor=south east, font=\small] {inflection $(0,\tfrac12)$};
\node[green!50!black] at (axis cs:-3.8,0.9) {concave up: $f''>0$};
\node[trapred] at (axis cs:3.8,0.25) {concave down: $f''<0$};
\end{axis}
\end{tikzpicture}
\caption{The logistic curve $f(x)=1/(1+e^{-x})$. For $x<0$, $f''>0$ (concave up,
curve above tangent); for $x>0$, $f''<0$ (concave down, curve below tangent).
At the inflection point $\big(0,\tfrac12\big)$, the tangent line (dashed)
slices directly through the curve, marking the transition from acceleration to deceleration.}
\label{ch16:fig-logistic}
\end{figure}

\begin{happeningbox}
What is the second derivative \emph{really} measuring at an inflection point?
Think of $f$ as position, so $f'$ is velocity and $f''$ is acceleration. Concave up
means accelerating (speeding up in the positive direction); concave down means
decelerating. The inflection point is the instant the acceleration passes through
zero and reverses --- the moment you stop pressing the gas and start pressing the
brake. The car is still moving forward (maybe even still gaining ground for a
while, because velocity has not yet hit zero), but the \emph{trend has turned}.
That is the deep content of $f''$ changing sign: not that motion stops, but that
the \emph{character} of the motion inverts. A peak in $f$ is where $f'=0$; an
inflection in $f$ is where $f'$ is at its own peak or valley. Inflection points of
$f$ are the extrema of $f'$.
\end{happeningbox}

\section{The second-derivative test: curvature as a shortcut}

The first-derivative test always works, but it asks you to examine $f'$ on
\emph{both} sides of a critical point. When $f''$ is easy to compute, there is a
faster route that looks at a single point. The idea is geometric and immediate:
at a critical point the tangent is horizontal, so the question ``peak or valley?''
becomes purely a question of \emph{which way the curve bends there}. Bends up
($\smile$) over a horizontal tangent: a valley. Bends down ($\frown$): a peak.

\begin{theorem}[Second-derivative test]\label{ch16:thm-sdt}
Suppose $f'(c)=0$ and $f''(c)$ exists.
\begin{itemize}
\item If $f''(c)>0$, then $f$ has a \emph{local minimum} at $c$.
\item If $f''(c)<0$, then $f$ has a \emph{local maximum} at $c$.
\item If $f''(c)=0$, the test is \emph{inconclusive} --- $c$ could be a max, a min,
or neither.
\end{itemize}
\end{theorem}

\begin{whybox}
Here is why the test works, tied straight back to concavity. Suppose
$f''(c)>0$. Recall (from the box above) that this means the curve lies \emph{above}
its tangent line at $c$ in a neighborhood. But the tangent line at $c$ is
\emph{horizontal}, because $f'(c)=0$ --- it is the constant line $y=f(c)$. So near
$c$ the curve lies above the horizontal line $y=f(c)$, i.e.\ $f(x)\ge f(c)$ nearby:
a local minimum. If $f''(c)<0$ the curve lies below its horizontal tangent, giving
$f(x)\le f(c)$: a local maximum. The second-derivative test is simply
``curve-above-tangent'' applied where the tangent happens to be horizontal.

A slightly more careful phrasing for $f''(c)>0$: by definition
$f''(c)=\lim_{x\to c}\frac{f'(x)-f'(c)}{x-c}>0$, and since $f'(c)=0$ this says
$\frac{f'(x)}{x-c}>0$ for $x$ near $c$. So $f'(x)$ and $(x-c)$ have the same sign:
$f'<0$ just left of $c$ and $f'>0$ just right. That is precisely the $-\to+$
pattern the first-derivative test reads as a minimum. The two tests agree because
the sign of the second derivative governs the sign change of the first.
\end{whybox}

\begin{example}[Second-derivative test, fast]\label{ch16:ex-sdt}
Reclassify the critical points of $f(x)=x^3-3x^2-9x+5$ using $f''$.

\smallskip
We have $f'(x)=3x^2-6x-9$, so $f''(x)=6x-6=6(x-1)$. Evaluate at the two critical
points:
\[
f''(-1)=6(-1-1)=-12<0 \ \Rightarrow\ \text{local max at }x=-1,
\]
\[
f''(3)=6(3-1)=12>0 \ \Rightarrow\ \text{local min at }x=3.
\]
Same verdict as the sign chart in Example~\ref{ch16:ex-fdt}, with no intervals to
test --- two evaluations and done. (Notice in passing that $f''=0$ at $x=1$, with
$f''$ flipping sign there: $x=1$ is an inflection point, the midpoint between the
peak and the valley, as cubics always have.)
\end{example}

\begin{trapbox}
When $f''(c)=0$, the second-derivative test tells you \emph{nothing} --- do not
conclude ``not an extremum.'' Compare three functions at $x=0$, all with $f'(0)=0$
and $f''(0)=0$:
\[
x^4 \ (\text{min}),\qquad -x^4 \ (\text{max}),\qquad x^3 \ (\text{neither}).
\]
The second derivative cannot tell these apart. When it vanishes, fall back on the
first-derivative test (check the sign of $f'$ on both sides) --- that one never
fails to give an answer.
\end{trapbox}

\begin{aestheticbox}
Why do we have two tests for the same thing? Because they answer different
practical questions. The first-derivative test is \emph{robust} --- it works at
corners and cusps where $f''$ does not even exist, and it is never inconclusive ---
but it costs you a study of $f'$ on a whole neighborhood. The second-derivative
test is \emph{cheap} --- one number per critical point --- but it demands a second
derivative and shrugs when that derivative is zero. Real problem-solving uses both:
reach for $f''$ first because it is fast, and when it returns zero, retreat to the
sign chart, which always closes the case. Knowing which tool to grab, and why each
can fail, is the mark of understanding rather than memorizing.
\end{aestheticbox}

\section{Reading $f$, $f'$, and $f''$ together: curve sketching}

Now everything assembles into a single skill: drawing a graph you have never
seen, accurately, from its formula alone. The discipline is to read three layers
and let them constrain one another --- three independent witnesses to the same
shape.

\begin{focusbox}
The three layers of a graph:
\begin{center}
\begin{tabular}{lll}
\toprule
\textbf{Object} & \textbf{What it controls} & \textbf{What you read from it}\\
\midrule
$f$ & height & intercepts, domain, asymptotes, symmetry\\
$f'$ & slope & increasing/decreasing, local max/min\\
$f''$ & bending & concave up/down, inflection points\\
\bottomrule
\end{tabular}
\end{center}
$f'$ governs whether the curve goes up or down. $f''$ governs whether it smiles or
frowns. Together they pin down the shape between your plotted points, which is
exactly the information naive dot-plotting was missing.
\end{focusbox}

A reliable checklist, in order:

\begin{enumerate}
\item \textbf{Domain and intercepts.} Where is $f$ defined? Where does it cross the
axes ($f=0$, and $f(0)$)?
\item \textbf{Symmetry and ends.} Even ($f(-x)=f(x)$, mirror across the $y$-axis)?
Odd ($f(-x)=-f(x)$, rotational)? What happens as $x\to\pm\infty$ --- any horizontal
or slant asymptotes (straight guide lines the curve approaches ever closer far out)? Any vertical asymptotes where $f$ blows up to $\pm\infty$ near an excluded point?
\item \textbf{First derivative.} Compute $f'$, find critical points, build the sign
chart: where increasing, where decreasing, and classify the extrema.
\item \textbf{Second derivative.} Compute $f''$, find where it is zero or
undefined, build its sign chart: concavity intervals and inflection points.
\item \textbf{Assemble.} Plot the special points (intercepts, extrema,
inflections, asymptotes), then connect them \emph{respecting} the slope and
concavity you found. There is now exactly one curve that fits all the constraints.
\end{enumerate}

\begin{example}[A full sketch: a quartic]\label{ch16:ex-sketch-quartic}
Sketch $f(x)=x^4-4x^3$.

\smallskip
\textbf{1. Domain, intercepts.} Defined for all $x$. Factor:
$f(x)=x^3(x-4)$, so $f=0$ at $x=0$ (a triple root --- the curve flattens through
the axis there) and at $x=4$. Also $f(0)=0$.

\textbf{2. Ends.} As $x\to\pm\infty$, the $x^4$ term dominates and $f\to+\infty$ on
both sides. No symmetry ($f(-x)=x^4+4x^3\neq\pm f(x)$). No asymptotes.

\textbf{3. First derivative.} $f'(x)=4x^3-12x^2=4x^2(x-3)$. Critical points where
$f'=0$: $x=0$ (double) and $x=3$. Sign chart:
\[
\begin{array}{c|ccc}
& (-\infty,0) & (0,3) & (3,\infty)\\\hline
4x^2 & + & + & +\\
(x-3) & - & - & +\\
f' & - & - & +\\
f & \searrow & \searrow & \nearrow
\end{array}
\]
At $x=3$, $f'$ goes $-\to+$: \textbf{local min}, $f(3)=81-108=-27$. At $x=0$, $f'$
does \emph{not} change sign ($-$ to $-$): \emph{not} an extremum, just a momentary
flat spot on the way down.

\textbf{4. Second derivative.} $f''(x)=12x^2-24x=12x(x-2)$. Zero at $x=0$ and
$x=2$. Sign chart:
\[
\begin{array}{c|ccc}
& (-\infty,0) & (0,2) & (2,\infty)\\\hline
12x & - & + & +\\
(x-2) & - & - & +\\
f'' & + & - & +\\
\text{shape} & \smile & \frown & \smile
\end{array}
\]
Concavity flips at \emph{both} $x=0$ and $x=2$ --- two inflection points:
$\big(0,0\big)$ and $\big(2,f(2)\big)=\big(2,-16\big)$.

\textbf{5. Assemble.} Coming from the upper left the curve falls, concave up, until
$x=0$ where it momentarily flattens (slope $0$) and changes to concave down ---
it keeps falling but now bulging downward --- through to $x=2$, where it switches
back to concave up while still descending, reaching the bottom of the valley at
$(3,-27)$, then climbing away to $+\infty$. The flat spot at the origin (slope zero
but \emph{not} an extremum) and the deep minimum at $x=3$ are the two features dot
plotting would most likely get wrong.
\end{example}

\begin{figure}[H]
\centering
\begin{tikzpicture}
\begin{axis}[
  width=12.5cm, height=6.8cm,
  axis lines=middle,
  xlabel={$x$}, ylabel={$y$},
  xmin=-1.6, xmax=4.8, ymin=-32, ymax=22,
  xtick={0,2,3,4}, ytick=\empty,
  domain=-1.4:4.4, samples=180, clip=false,
]
\addplot[curveblue, very thick] {x^4 - 4*x^3};
\addplot[only marks, mark=*, mark size=1.6pt, tangentred]
  coordinates {(3,-27)};
\addplot[only marks, mark=*, mark size=1.6pt, structurepurple]
  coordinates {(0,0) (2,-16)};
\node[anchor=north, tangentred] at (axis cs:3,-28.5) {min $(3,-27)$};
\node[anchor=south west, structurepurple] at (axis cs:0.1,0.5) {infl.\ $(0,0)$};
\node[anchor=west, structurepurple] at (axis cs:2.1,-16) {infl.\ $(2,-16)$};
\node[anchor=south east, curveblue] at (axis cs:4.1,16) {$f$};
\end{axis}
\end{tikzpicture}
\caption{$f(x)=x^4-4x^3$. A flat-but-not-extreme spot at the origin, a genuine
minimum at $(3,-27)$, and two inflection points (purple) at $x=0$ and $x=2$ where
the bending reverses.}
\label{ch16:fig-quartic}
\end{figure}

\begin{example}[A sketch with asymptotes: a rational function]\label{ch16:ex-sketch-rational}
Sketch $f(x)=\dfrac{x^2}{x^2-1}$.

\smallskip
\textbf{1. Domain, intercepts.} The denominator $x^2-1=(x-1)(x+1)$ vanishes at
$x=\pm1$, so those are excluded: vertical asymptotes there. $f=0$ only when the
numerator $x^2=0$, i.e.\ $x=0$, where also $f(0)=0$. So the only intercept is the
origin.

\textbf{2. Symmetry, ends.} $f(-x)=\dfrac{x^2}{x^2-1}=f(x)$: \emph{even}, symmetric
across the $y$-axis. As $x\to\pm\infty$, $f=\dfrac{x^2}{x^2-1}\to 1$: a horizontal
asymptote $y=1$. Near the vertical asymptotes the sign of the denominator tells the
story: for $\abs{x}$ slightly less than $1$, $x^2-1<0$ so $f\to-\infty$; for
$\abs{x}$ slightly more than $1$, $x^2-1>0$ so $f\to+\infty$.

\textbf{3. First derivative.} By the quotient rule,
\[
f'(x)=\frac{(2x)(x^2-1)-x^2(2x)}{(x^2-1)^2}
=\frac{2x^3-2x-2x^3}{(x^2-1)^2}
=\frac{-2x}{(x^2-1)^2}.
\]
The denominator is positive wherever defined, so the sign of $f'$ is the sign of
$-2x$: positive for $x<0$, negative for $x>0$. The only critical point is $x=0$,
and $f'$ goes $+\to-$ there: a \textbf{local maximum} at $(0,0)$. (It is a maximum
of the central branch, even though outer branches take much larger values,
illustrating that ``local'' refers strictly to an open neighborhood.)

\textbf{4. Second derivative.} Differentiating $f'(x)=\dfrac{-2x}{(x^2-1)^2}$ with
the quotient rule:
\[
f''(x)=\frac{-2(x^2-1)^2 - (-2x)\cdot 2(x^2-1)(2x)}{(x^2-1)^4}
=\frac{-2(x^2-1) + 8x^2}{(x^2-1)^3}
=\frac{2(3x^2+1)}{(x^2-1)^3}.
\]
The numerator $2(3x^2+1)$ is always positive, so the sign of $f''$ matches the sign
of $(x^2-1)^3$, which matches the sign of $(x^2-1)$:
\[
f''>0 \text{ when } \abs{x}>1 \ (\text{concave up}),\qquad
f''<0 \text{ when } \abs{x}<1 \ (\text{concave down}).
\]
There is no inflection point, because the sign change of $f''$ happens only at
$x=\pm1$, where $f$ is not even defined.

\textbf{5. Assemble.} Three branches. The middle branch (between the asymptotes)
is a downward-bulging arch peaking at the origin and diving to $-\infty$ at both
$x=\pm1$. The two outer branches descend from $+\infty$ (just outside the
asymptotes) toward the horizontal line $y=1$, concave up, mirror images of each
other.
\end{example}

\begin{figure}[H]
\centering
\begin{tikzpicture}
\begin{axis}[
  width=12.5cm, height=7.2cm,
  axis lines=middle,
  xlabel={$x$}, ylabel={$y$},
  xmin=-4.2, xmax=4.2, ymin=-5.2, ymax=5.2,
  xtick=\empty, ytick=\empty,
  clip=true,
  restrict y to domain=-6:6,
]
\addplot[curveblue, very thick, domain=-4:-1.05, samples=120] {x^2/(x^2-1)};
\addplot[curveblue, very thick, domain=-0.96:0.96, samples=120] {x^2/(x^2-1)};
\addplot[curveblue, very thick, domain=1.05:4, samples=120] {x^2/(x^2-1)};
\addplot[gray, dashed, thick] coordinates {(-1,-5.2) (-1,5.2)};
\addplot[gray, dashed, thick] coordinates {(1,-5.2) (1,5.2)};
\addplot[gray, dashed, thick, domain=-4:4] {1};
\addplot[only marks, mark=*, mark size=1.6pt, tangentred] coordinates {(0,0)};
\node[anchor=south, tangentred, fill=white, inner sep=1.2pt, font=\footnotesize] at (axis cs:0,0.22) {max $(0,0)$};
\node[anchor=north east, gray!70!black, font=\footnotesize] at (axis cs:-1.08,4.4) {$x=-1$};
\node[anchor=north west, gray!70!black, font=\footnotesize] at (axis cs:1.08,4.4) {$x=1$};
\node[anchor=south, gray!70!black, font=\footnotesize] at (axis cs:3.2,1.15) {$y=1$};
\end{axis}
\end{tikzpicture}
\caption{$f(x)=x^2/(x^2-1)$. Vertical asymptotes at $x=\pm1$ (dashed) split the
graph into three branches; horizontal asymptote $y=1$. The central arch has a local
maximum at the origin.}
\label{ch16:fig-rational}
\end{figure}

\begin{walkbox}
Here is a graph you can build entirely in your head on a walk. Imagine a function
whose \emph{derivative} $f'$ looks like a single hump --- zero at $x=0$, rising to a
peak at $x=2$, falling back to zero at $x=4$, and staying near zero outside. Now ask
yourself, without paper: where is $f$ itself increasing? (Wherever $f'>0$, which is
the whole stretch from $0$ to $4$ --- so $f$ climbs the entire time.) Where is $f$
\emph{concave up}? (Wherever $f'$ is \emph{increasing}, i.e.\ from $0$ to $2$.)
Where is $f$'s inflection point? (At $x=2$, the peak of $f'$ --- the steepest point
of the climb.) You have just read the shape of $f$ off the shape of $f'$, which is
the single most useful fluency in this chapter (Figure~\ref{ch16:fig-f-and-fprime}
shows a worked instance). Practice going both directions: from
$f$ to $f'$, and from $f'$ back up to $f$.
\end{walkbox}

\begin{figure}[H]
\centering
\begin{tikzpicture}
\begin{axis}[
  name=laneF,
  width=11.5cm, height=3.8cm,
  axis lines=middle,
  ylabel={$f(x)$},
  xmin=-2.4, xmax=2.4, ymin=-2.8, ymax=2.8,
  xtick=\empty, ytick=\empty,
  domain=-2.1:2.1, samples=120, clip=false,
  title style={font=\small},
]
\addplot[curveblue, very thick] {x^3 - 3*x};
\addplot[only marks, mark=*, mark size=1.6pt, tangentred] coordinates {(-1,2) (1,-2)};
\addplot[only marks, mark=*, mark size=1.6pt, structurepurple] coordinates {(0,0)};
\node[anchor=south, tangentred] at (axis cs:-1,2.1) {\footnotesize max};
\node[anchor=north, tangentred] at (axis cs:1,-2.1) {\footnotesize min};
\node[anchor=south west, structurepurple] at (axis cs:0.08,0.1) {\footnotesize inflection};
\node[anchor=center, curveblue] at (axis cs:1.45,1.0) {\footnotesize $f=x^3-3x$};
\coordinate (laneF_m1) at (axis cs:-1,2.6);
\coordinate (laneF_0)  at (axis cs:0,2.6);
\coordinate (laneF_p1) at (axis cs:1,2.6);
\end{axis}

\begin{axis}[
  at={(laneF.below south west)}, anchor=north west, yshift=-0.35cm,
  name=laneFp,
  width=11.5cm, height=3.8cm,
  axis lines=middle,
  ylabel={$f'(x)$},
  xmin=-2.4, xmax=2.4, ymin=-3.6, ymax=6.5,
  xtick=\empty, ytick=\empty,
  domain=-2.1:2.1, samples=120, clip=false,
  title style={font=\small},
]
\addplot[tangentred, very thick] {3*x^2 - 3};
\addplot[only marks, mark=*, mark size=1.6pt, tangentred] coordinates {(-1,0) (1,0)};
\node[green!50!black] at (axis cs:-1.7,2.5) {\footnotesize $f'>0$};
\node[trapred] at (axis cs:-0.5,-1.3) {\footnotesize $f'<0$};
\node[green!50!black] at (axis cs:1.7,2.5) {\footnotesize $f'>0$};
\node[anchor=center, tangentred] at (axis cs:-0.5,3.8) {\footnotesize $f'=3x^2-3$};
\end{axis}

\begin{axis}[
  at={(laneFp.below south west)}, anchor=north west, yshift=-0.35cm,
  name=laneFpp,
  width=11.5cm, height=3.4cm,
  axis lines=middle,
  xlabel={$x$}, ylabel={$f''(x)$},
  xmin=-2.4, xmax=2.4, ymin=-7.5, ymax=7.5,
  xtick=\empty, ytick=\empty,
  domain=-2.1:2.1, samples=60, clip=false,
  title style={font=\small},
]
\addplot[structurepurple, very thick] {6*x};
\addplot[only marks, mark=*, mark size=1.6pt, structurepurple] coordinates {(0,0)};
\node[anchor=north, trapred] at (axis cs:-1.6,-2.5) {\footnotesize $f''<0$ ($\frown$)};
\node[anchor=south, green!50!black] at (axis cs:1.6,2.5) {\footnotesize $f''>0$ ($\smile$)};
\node[anchor=center, structurepurple] at (axis cs:-0.5,4.0) {$f''=6x$};
\coordinate (laneFpp_m1) at (axis cs:-1,-7.2);
\coordinate (laneFpp_0)  at (axis cs:0,-7.2);
\coordinate (laneFpp_p1) at (axis cs:1,-7.2);
\end{axis}

\draw[gray!50, dashed, thick] (laneF_m1) -- (laneFpp_m1);
\draw[structurepurple!50, dashed, thick] (laneF_0) -- (laneFpp_0);
\draw[gray!50, dashed, thick] (laneF_p1) -- (laneFpp_p1);
\node[anchor=north, gray!70!black] at (laneFpp_m1) {\footnotesize $x=-1$};
\node[anchor=north, structurepurple] at (laneFpp_0) {\footnotesize $x=0$};
\node[anchor=north, gray!70!black] at (laneFpp_p1) {\footnotesize $x=1$};
\end{tikzpicture}
\caption{The three synchronized layers of curve analysis for $f(x)=x^3-3x$.
Top lane: the graph of $f$ rises where $f'>0$ and falls where $f'<0$.
Middle lane: the zeros of $f'$ at $x=\pm1$ pinpoint the peak and valley of $f$.
Bottom lane: the zero of $f''$ at $x=0$ marks the inflection point, where $f'$ attains
its global minimum (steepest downhill descent) and the concavity of $f$ flips from
down ($\frown$) to up ($\smile$).}
\label{ch16:fig-f-and-fprime}
\end{figure}

\begin{labbox}
Take $f(x)=x e^{-x}$ --- a shape that shows up constantly in physics (radioactive
decay times a linear rise), pharmacology (drug concentration over time), and
probability. Work out the shape entirely by hand, then check your sketch against a
plot.
\begin{enumerate}
\item Show $f'(x)=e^{-x}(1-x)$. Find the one critical point and classify it with
the first-derivative test. (Answer: max at $x=1$, value $1/e\approx0.368$.)
\item Show $f''(x)=e^{-x}(x-2)$. Find the inflection point. (Answer: at $x=2$,
where the decay's concavity flips from down to up.)
\item As $x\to+\infty$, does $f\to0$ or $\to\infty$? (The exponential
$e^{-x}$ beats the linear $x$, so $f\to0^+$ --- a horizontal asymptote $y=0$.) As
$x\to-\infty$? ($f\to-\infty$, because $e^{-x}$ explodes.)
\item Assemble: a curve rising from below, crossing the origin, climbing to a single
peak at $x=1$, then decaying back toward zero, with the decay easing off (inflection)
at $x=2$. Sketch it, then plot it and compare.
\end{enumerate}
\end{labbox}

\begin{connectionbox}
Concavity is the one-variable shadow of a much bigger idea: \emph{convexity}, which
runs through optimization, economics, and machine learning. A function with
$f''\ge0$ everywhere is called \emph{convex}; the curve-above-its-chords picture is
its defining property. Convex functions have a key property you can already
anticipate from this chapter: \emph{any local minimum is automatically a global
minimum} --- there are no isolated suboptimal valleys, because the curve never
turns back down once it has started up. This is why training a model whose
loss function is convex is tractable (gradient descent cannot be trapped) and why
non-convex losses with multiple local minima present greater challenges. The
second-derivative test is the one-variable origin of that theory. We will meet the
multivariable version --- where $f''$ becomes a matrix and ``concave up'' becomes a
statement about that matrix --- when we reach functions of several variables.
\end{connectionbox}

\section{Finding the global extremes on an interval}

Local extrema are about neighborhoods. But the questions that matter in practice
are usually \emph{global}: over this whole range of choices, what is the absolute
best? On an unbounded domain a function may have no global max at all (think
$f(x)=x$, climbing forever). The clean, decidable case is a function on a
\emph{closed, bounded} interval $[a,b]$, and there the Extreme Value Theorem
(from the continuity chapter) guarantees both a global max and a global min
actually exist. The only question is \emph{where}.

\begin{focusbox}
\textbf{The closed-interval method.} To find the absolute max and min of a
continuous $f$ on $[a,b]$:
\begin{enumerate}
\item Find all critical points of $f$ \emph{inside} $(a,b)$.
\item Make a list: those critical points, plus the two endpoints $a$ and $b$.
\item Evaluate $f$ at every point on the list.
\item The largest value is the global max; the smallest is the global min.
\end{enumerate}
No sign charts, no second derivatives --- just evaluate and compare. The logic is
airtight: a global extremum is in particular a local one \emph{unless} it sits at
an endpoint (where ``local'' has no left or right neighbor to compare with). Local
interior extrema live at critical points (Fermat); the only other place is an
endpoint. So the global champion is somewhere on the list, guaranteed.
\end{focusbox}

\begin{example}[Closed-interval optimization]\label{ch16:ex-closed}
Find the absolute maximum and minimum of $f(x)=x^3-3x^2-9x+5$ on $[-2,4]$.

\smallskip
We already found the critical points $x=-1$ and $x=3$ (both inside $(-2,4)$). The
list is $\{-2,-1,3,4\}$. Evaluate:
\[
\begin{aligned}
f(-2)&=-8-12+18+5=3,\\
f(-1)&=-1-3+9+5=10,\\
f(3)&=27-27-27+5=-22,\\
f(4)&=64-48-36+5=-15.
\end{aligned}
\]
The largest is $f(-1)=10$: \textbf{global maximum}. The smallest is $f(3)=-22$:
\textbf{global minimum}. Notice that the endpoint $x=4$ gives $-15$, lower than the
\emph{other} endpoint and lower than the left critical point --- a reminder that
endpoints are genuine contenders and must always be on the list. Had we forgotten
them, we would still have caught the true extremes here, but on a different interval
an endpoint could easily be the winner (Figure~\ref{ch16:fig-closed-interval}).
\end{example}

\begin{figure}[H]
\centering
\begin{tikzpicture}
\begin{axis}[
  width=12cm, height=6.6cm,
  axis lines=middle,
  xlabel={$x$}, ylabel={$y$},
  xmin=-2.8, xmax=4.8, ymin=-26, ymax=16,
  xtick=\empty, ytick=\empty,
  domain=-2:4, samples=120, clip=false,
]
\addplot[curveblue, very thick] {x^3 - 3*x^2 - 9*x + 5};
\addplot[gray!60!black, dashed, thick] coordinates {(-2,0) (-2,3)};
\addplot[gray!60!black, dashed, thick] coordinates {(-1,0) (-1,10)};
\addplot[gray!60!black, dashed, thick] coordinates {(3,0) (3,-22)};
\addplot[gray!60!black, dashed, thick] coordinates {(4,0) (4,-15)};
\addplot[only marks, mark=*, mark size=2pt, sourceblue]
  coordinates {(-2,3) (4,-15)};
\addplot[only marks, mark=*, mark size=2.2pt, tangentred]
  coordinates {(-1,10) (3,-22)};
\node[anchor=south east, sourceblue, font=\small] at (axis cs:-2,3.3) {endpoint $(-2,3)$};
\node[anchor=south, tangentred, font=\small, align=center] at (axis cs:-1.1,10.6) {\textbf{global max}\\ $(-1,10)$};
\node[anchor=north, tangentred, font=\small, align=center] at (axis cs:3,-22.8) {\textbf{global min}\\ $(3,-22)$};
\node[anchor=west, sourceblue, font=\small] at (axis cs:4.05,-15) {endpoint $(4,-15)$};
\end{axis}
\end{tikzpicture}
\caption{The closed-interval method on $[-2,4]$ for $f(x)=x^3-3x^2-9x+5$. Four
candidates compete: the two endpoints and the two interior critical points.
Evaluating all four reveals the global maximum at $(-1,10)$ and the global minimum
at $(3,-22)$. Notice that the right endpoint $(4,-15)$ sits lower than the left
endpoint, confirming that endpoints are genuine contenders.}
\label{ch16:fig-closed-interval}
\end{figure}

\begin{trapbox}
The closed-interval method needs both hypotheses of the Extreme Value Theorem:
\emph{continuous} function on a \emph{closed, bounded} interval. Drop either and the
guarantee evaporates. On the open interval $(0,1)$ the function $f(x)=x$ has no
maximum --- it gets arbitrarily close to $1$ but never reaches it, and $1$ is not in
the domain. On $[0,\infty)$ the function $f(x)=x$ has no maximum either. And a
function with a vertical asymptote inside the interval (a discontinuity) can shoot
to $+\infty$ with no max at all. When the EVT does not apply, finding global
extrema is a genuinely subtler business --- you must examine the behavior at the
open ends or at infinity by hand.
\end{trapbox}

\section{What we built}

The whole chapter is one idea unfolding: \emph{the derivatives of $f$ encode the
shape of $f$}. The first derivative records direction --- its zeros and
nondefined points (critical points) trap every local extremum, and its sign on
either side (the first-derivative test) classifies them. The second derivative
records bending --- its sign gives concavity, its sign \emph{changes} give
inflection points, and at a flat tangent its sign settles peak-versus-valley in
one stroke (the second-derivative test). Reading both at once, you can reconstruct
a graph exactly, and on a closed interval you can pin down the absolute best and
worst with only a short list of values.

Underneath every one of these tests sits the Mean Value Theorem, quietly
converting facts about $f'$ into facts about $f$. ``Positive derivative means
increasing'' is the MVT; the first-derivative test is that fact applied on each
side of a critical point; concavity is that same fact applied to $f'$; the
second-derivative test is concavity applied at a horizontal tangent. Those four rules are really one theorem doing four jobs.

In the next chapter we put this machine to work on real questions --- the largest
box, the cheapest can, the shortest path, the maximum profit --- where the art is
no longer running the tests but \emph{setting up} the right function to optimize.
The hard part of optimization was never the calculus. It was deciding what to take
the derivative of.

\subsection*{Summary: Curve Shape and Optimization}
\addcontentsline{toc}{subsection}{Summary: Curve Shape and Optimization}

\begin{connectionbox}
Armed with the first and second derivatives, you can systematically determine:
\begin{itemize}
  \item \emph{Exact critical points and extrema}: Locating every candidate peak, valley, and plateau ($f'=0$ or undefined) without plotting millions of points.
  \item \emph{Global monotonicity and concavity}: Partitioning the real line into intervals where the graph rises, falls, bends upward like a cup, or sheds water like a dome.
  \item \emph{Complete shape and curvature}: Identifying inflections where physical quantities change acceleration, and using the closed-interval method to guarantee global optima.
\end{itemize}

Yet the critical-point framework relies on isolated critical points and well-behaved, continuous derivatives. It falters on functions with infinite-frequency oscillations. Consider the canonical counterexample:
\[
  f(x) = \begin{cases} x^2 \sin\big(\tfrac{1}{x}\big), & x \ne 0, \\ 0, & x = 0. \end{cases}
\]
Here $f'(0) = \lim_{h\to 0} \frac{h^2 \sin(1/h)}{h} = 0$, so $x=0$ is a critical point. Yet for $x \ne 0$, $f'(x) = 2x\sin(1/x) - \cos(1/x)$, which oscillates between positive and negative infinitely often in every open neighborhood of $0$. There are infinitely many local maxima and minima accumulating at the origin! The first derivative test cannot be applied because $f'$ does not have constant sign on either side of $0$, and $f''(0)$ does not even exist. Standard calculus tools assume the tame, well-behaved derivatives typical of physical applications.

In Chapter~17 (\emph{Optimization: Finding the Best}), we take these classification tools into the physical world. We will solve Fermat's principle of least time, derive Snell's law of optical refraction, and solve structural and economic design problems by translating real-world constraints into single-variable extrema.
\end{connectionbox}

\section*{Exercises}
\addcontentsline{toc}{section}{Exercises}

\begin{exercise}[Mental warm-ups]\label{ch16:ex-warmup}
Answer each question mentally or with a one-line sketch before calculating:
\begin{enumerate}
\item If $f'(x) = (x-2)^2(x-5)$, where does $f$ have a local minimum? Does $x=2$ yield an extremum?
\item Can a cubic polynomial $f(x)=ax^3+bx^2+cx+d$ ($a\ne0$) have two local maxima? Explain using the degree of $f'$.
\item If a differentiable function on $[1,5]$ satisfies $f'(x)>0$ for all $x\in(1,5)$, where do its global minimum and global maximum occur?
\end{enumerate}
\end{exercise}

\begin{exercise}\label{ch16:exA}
For each function, find all critical points and classify each as a local maximum,
local minimum, or neither, using the first-derivative test.
\begin{enumerate}
\item $f(x)=2x^3-9x^2+12x-3$
\item $f(x)=x^4-8x^2+3$
\item $f(x)=x^{2/3}(x-5)$ \quad(watch for a point where $f'$ does not exist)
\end{enumerate}
\end{exercise}

\begin{exercise}\label{ch16:exB}
Find the intervals of concavity and all inflection points of
$f(x)=x^4-6x^2+5$. Then confirm, using the second-derivative test, the nature of
each critical point.
\end{exercise}

\begin{exercise}\label{ch16:exC}
Find the absolute maximum and minimum of $f(x)=2x^3-3x^2-12x+1$ on the closed
interval $[-2,3]$. Identify whether each extreme occurs at an endpoint or at an
interior critical point.
\end{exercise}

\begin{exercise}\label{ch16:exD}
Carry out a full curve sketch of $f(x)=\dfrac{x}{x^2+1}$: domain, symmetry,
intercepts, end behavior and asymptotes, intervals of increase/decrease with
extrema, and intervals of concavity with inflection points. Then draw the graph.
(Hint: $f$ is odd, $f\to0$ at both ends, and you will find one max, one min, and
three inflection points.)
\end{exercise}

\begin{exercise}\label{ch16:exE}
\textbf{(Conceptual.)} For each description, either sketch a function meeting it or
explain briefly why none can exist.
\begin{enumerate}
\item $f'(2)=0$ but $f$ has neither a max nor a min at $x=2$.
\item $f''(2)=0$ but $f$ has no inflection point at $x=2$.
\item $f$ has a local minimum at $x=2$ but $f'(2)$ does not exist.
\item $f$ has a local maximum at $x=2$ and $f''(2)>0$.
\end{enumerate}
\end{exercise}

\begin{exercise}\label{ch16:exF}
\textbf{(Reading $f$ from $f'$.)} The derivative of a function $f$ is
$f'(x)=(x-1)^2(x-4)$. Without finding $f$ itself, determine: (a) the intervals
where $f$ is increasing and decreasing; (b) the locations and types of all local
extrema; (c) whether $x=1$ is special, and why the squared factor matters.
\end{exercise}

\begin{exercise}\label{ch16:exG}
\textbf{(Why the test works.)} Suppose $f''$ is continuous and $f''(c)>0$ at a point
where $f'(c)=0$. Using the definition $f''(c)=\lim_{x\to
c}\dfrac{f'(x)-f'(c)}{x-c}$, prove directly that $f'$ is negative just to the left
of $c$ and positive just to the right, and conclude that $f$ has a local minimum at
$c$. (This is the second-derivative test, derived from scratch.)
\end{exercise}

\begin{exercise}\label{ch16:exH}
\textbf{(Proof, tangent line and concavity.)} Let $f$ be twice differentiable with
$f''(x)>0$ for all $x$ in an interval $I$. Prove that for any fixed $a\in I$,
\[
f(x)\ \ge\ f(a)+f'(a)(x-a)\qquad\text{for all } x\in I,
\]
with equality only at $x=a$. (That is: a strictly convex curve lies strictly above
each of its tangent lines. Imitate the gap-function argument from the chapter:
let $g(x)=f(x)-f(a)-f'(a)(x-a)$ and study the sign of $g'$.)
\end{exercise}

\begin{exercise}\label{ch16:exI}
\textbf{(Inflections are extrema of $f'$.)} Explain carefully why an inflection
point of $f$ corresponds to a local maximum or minimum of $f'$. Use this to find
the inflection point of $f(x)=\dfrac{1}{1+e^{-x}}$ by maximizing $f'$ instead of
analyzing $f''$ directly, and confirm you get $x=0$.
\end{exercise}

\begin{exercise}\label{ch16:exJ}
\textbf{(Modeling.)} A drug's concentration in the bloodstream $t$ hours after a
dose is modeled by $C(t)=12t\,e^{-t/2}$ for $t\ge0$. (a) When is the concentration
highest, and what is that peak concentration? (b) The concentration is ``falling
fastest'' at the inflection point of $C$; find that time. (c) Sketch $C$ for
$t\in[0,12]$, marking the peak and the inflection.
\end{exercise}
